Distributed inference at the edge is constrained by uplink bandwidth.
Requests that share a context can reuse key-value cache, which reduces both computation and transfer.
However, reuse couples request placement: a request that arrives at a node without the cache either rebuilds it or fetches it.
We study admission control when reuse is possible but not guaranteed.
Our working hypothesis is that a modest reservation for cache-bearing nodes improves tail latency without stranding capacity.
This hypothesis matters because operators already reserve capacity for warm caches, but they do so with fixed thresholds.
A fixed threshold cannot react to changes in request mix or to congestion on the uplink.
We therefore derive a threshold from observed reuse and transfer cost rather than from a static policy.
The resulting rule is simple enough to deploy and explicit enough to audit.
We evaluate it on a synthetic trace and report the conditions under which it helps.
We do not claim that the rule is optimal, and we identify the cases where it is not.
